% Introduction (DECISIONS §D1, §A11, §F29, §F30; framing: REFRAME.md, which deletes the paragraph "Relation to the author's
% notes and to other work" and moves the credit to one footnote and the acknowledgements). Sources: the body sections,
% research/prior/hanni-grokking-a-notion.md (no quotation). The former "Prior work" paragraph is in related.tex.
% The F29 sentence for the teacher's-mistakes bullet is the second half of the informal theorem below.
% Binding notes: the absorption objection is not attributed to a person in the running text; the single footnote crediting
% the proposal (REFRAME §4, at its first statement below) also records the online discussion.
\section{Introduction}\label{sec:intro}

Solomonoff function induction (SFI) is Bayesian prediction of a label $y$ from an input $x$ with a universal prior over random partial functions \citep{claudeSFI2026fundamentals}. It is the supervised-learning counterpart of Solomonoff induction (SI), and like SI it learns every computable regularity of the labels, including a labeller's systematic mistakes. A human learner, by contrast, can grasp what a teacher means by a word without learning to predict what the teacher would say, mistakes included. We study the following proposal for this purpose, a stronger simplicity prior: when the mistakes are explainable deviations from a simple model, select the simple model with noise, by penalising the length of a hypothesis more than its likelihood term.\footnote{The proposal first appeared in the author's note \citep{hanni2026grokking} and in an online discussion (on LessWrong), where the absorption objection described below was also raised. So did the suggestion that private randomness at each input blocks this objection, the three-hypothesis calculation below (in both versions of \cref{rem:teacher:hanni}), the question whether its good model is a vertex of the convex hull, and an absolute complexity bound as an alternative. SFI itself first appeared in the author's working notes \citep{hanni2026notes}.} In the form studied here, a hypothesis is a program $\pi$ of length $\ell(\pi)$ with likelihood $P_\pi$, and one minimises
\[
J_\lambda(\pi;D)=\lambda\,\ell(\pi)-\log P_\pi(D),\qquad\lambda>1,
\]
or uses the power prior $2^{-\lambda\ell(\pi)}$ (\cref{def:setup:objectives}). We consider it for function induction with private randomness at each input: a hypothesis is a probabilistic program that uses fresh coins at every input.

A standard objection, which we call the \emph{absorption objection}, is that under a steeper penalty a hypothesis that is itself a universal mixture is short, penalises complexity inside itself only in the usual way, and therefore absorbs all the structure of the data, errors included. In sequence prediction the objection holds (\cref{sec:absorb}). In function induction one may hope that it fails if each input samples its own random bits, so that a hypothesis pays its likelihood term separately at every input. A three-hypothesis calculation illustrates the proposal. Take a $100$-bit simple model and systematic errors on a $2^{-10}$ fraction of the inputs, whose set takes $10{,}000$ bits to specify; of three hypotheses (the simple model with noise, a model that learns the errors, and a fair coin) the first is selected for a range of $\lambda/n$ (\cref{rem:teacher:hanni}). The example is favourable to the proposal, and three hypotheses are not all hypotheses: does another hypothesis do better, and is some other pathology present? The split between model and noise arises in its simplest form for labels drawn by a simple sampler, whose coins are the noise; programs that fill in the sampled bits store the noise as model. At $\lambda=1$ the two get similar posterior weight, in SFI and in SI (\cref{thm:split:lambda1,prop:split:si}); does a steeper prior prefer the sampler? \Cref{tab:intro:claims} lists these and further questions and where they are answered.

\paragraph{Conventions.} Logarithms are to base $2$. $n$ is the number of samples, $N=2^d$, $c=\lambda/n$, and $\kappa$ is a constant that depends only on the reference machine and fixed programs; numbers computed with $\kappa=0$ are illustrations, not results about a reference machine. A shared hypothesis $P$ (\cref{def:setup:shared}) gives the label $y$ at a new input $x$ after data $D$ the probability $P(D+(x,y))/P(D)$; in this \emph{semimeasure convention} the cumulative log loss telescopes, in every order, to at most $-\log P(D)$ bits. In the \emph{normalised convention} the denominator is $\sum_{y'}P(D+(x,y'))\le P(D)$, so losses are never larger.

\paragraph{The main result.} \Cref{thm:teacher:main} makes the example rigorous for error rules that are incompressible relative to the model. Informally:

\smallskip\noindent\textbf{Theorem} (informal). \emph{Let the inputs lie in $\{0,1\}^d$, let the simple model be a random parity $f_a$, and let a teacher flip its labels on the inputs whose first $m$ bits lie in an error rule $S\subseteq\{0,1\}^m$ of density $p<0.056$, so that the error set $S\times\{0,1\}^{d-m}$ has product form. Score each hypothesis $\pi$ with private randomness by the population objective $c\,\ell(\pi)+R(\pi)$, with $R(\pi)$ its expected log loss on the teacher's labels. If $S$ is incompressible relative to the parity, the inputs are uniform and $\lambda=cn$ with $c$ in an explicit window $(2^{-m},c_{\rm hi})$, every near-optimal product hypothesis of the population objective knows the parity, is short, holds little information about $S$ and gives the erroneous labels average probability at most $0.199$ $(\kappa=0)$ or $0.274$ $(\kappa=30)$ at $c=10^{-4}$ and the example's sizes, against $1$ for a memoriser. Below about $2^{-m}$ the memoriser beats every hypothesis that knows the parity and fits no better than noise.}

\smallskip\noindent The example's sizes are $d=100$, $m=20$, $p=2^{-10}$ and $11{,}710$ bits for $S$; the numbers assume $\delta_S=0$ for the deficiency of $S$. There the window is $(9.54\cdot10^{-7},8.06\cdot10^{-3}]$ at $\kappa=0$ and $(9.54\cdot10^{-7},3.66\cdot10^{-3}]$ at $\kappa=30$ (\cref{rem:teacher:numbers}).

\paragraph{Contributions.}
\begin{itemize}[leftmargin=1.5em,itemsep=1pt]
\item \emph{Absorption} (\cref{sec:absorb}). In every class that contains its own Bayes mixture, such as SI and SFI with shared randomness, every $\lambda$-MAP is shorter than a bound that does not depend on the data, and for every fixed $\lambda>1$ a computable deterministic teacher's labels, errors included, are eventually predicted (\cref{thm:absorb:absorption}). No hypothesis with private randomness is universal (\cref{cor:absorb:nouniversal}); this protection costs a factor $\min(\lambda,2)$ on shared latent coins and ends with labelled history (\cref{prop:absorb:price,prop:absorb:history}).
\item \emph{The teacher's mistakes} (\cref{sec:teacher}). The theorem above holds over all product hypotheses (\cref{thm:teacher:main}), with versions for full truth tables and for sampled data (\cref{cor:teacher:fulltable,cor:teacher:sampled}).
\item \emph{Limits} (\cref{sec:limits}). Plain Bayes, MDL and SI learn the errors, and so does every schedule with $(\lambda_n/n)(d+L_S)\to0$ along a sequence of instances, so steeper selection needs $\lambda=\Theta(n)$ (\cref{prop:limits:plain,rem:limits:schedule}). For an error rule with the example's density and about its description length, no hypothesis that fits no better than the good model is near-optimal on almost a third ($\kappa=0$; a quarter at $\kappa=30$) of the window's logarithmic width (\cref{prop:limits:envelope}). Cheap, frequent error rules defeat every $c$ once the simple model is long enough (\cref{prop:limits:rateblind,lem:limits:robust}).
\item \emph{The model--randomness split} (\cref{sec:split}). At $\lambda=1$ the posterior odds of the fill-ins of a sampler of length $k$ against the sampler lie in $[2^{-2\log(t_n+1)-\kappa},2^{k+\delta^*}]$, with $2\log(t_n+1)$ doubly logarithmic in the input length and at most triply logarithmic in $n$ and in the sampled randomness, and $\delta^*$ a Ville deficiency (\cref{thm:split:lambda1}). For fixed $\lambda>1$ with private randomness no $\lambda$-MAP is a fill-in once the sampled randomness $-\log\nu(D_n)$ exceeds $\lambda(k+\delta^*)/(\lambda-1)$ (\cref{thm:split:private}); with shared randomness, and in SI, no sampler longer than a data-independent bound $\ell^*_\lambda$ is selected (\cref{thm:split:shared}).
\item \emph{Algorithmic statistics and computable classes} (\cref{sec:algstat,sec:computable}). For a noisy random parity, product sophistication exceeds set sophistication by $d-O(\log d)$ bits for a range of sufficiency slacks (\cref{prop:algstat:parity}). The lower bounds and the private-randomness split hold in every c.e.\ prefix-free product class (for the split, one that contains the sampler), with the class's own constants (\cref{prop:related:computable}).
\end{itemize}

\begin{table}[t]
\centering\footnotesize\crefabbrev\setlength{\tabcolsep}{4pt}
\begin{tabularx}{\linewidth}{@{}>{\raggedright\arraybackslash}p{4.6cm}>{\raggedright\arraybackslash}X>{\raggedright\arraybackslash}p{2.2cm}@{}}
\toprule
Question & Answer & Where\\
\midrule
Does absorption occur in the sequence case, and not if each input samples its own random bits? & Yes, in every class that contains its own mixture; and no hypothesis with private randomness is universal. In part: absorption returns with labelled history, and shared latent causes cannot be expressed & \cref{thm:absorb:absorption,cor:absorb:nouniversal,prop:absorb:history,prop:absorb:price}\\
Can the good model of the example be selected for a range of $c$, as a vertex of the convex hull of (complexity, expected log loss)? & Yes for error rules incompressible relative to the model. Not for every rule with the same density and description length: for one such rule, no hypothesis that fits no better than the good model is near-optimal on about a third ($\kappa=0$) to a quarter ($\kappa=30$) of the window's logarithmic width; elsewhere only the weaker bounds of \cref{thm:teacher:main} with that rule's deficiency apply & \cref{thm:teacher:main,prop:limits:envelope}\\
Does another hypothesis do better in the example? & No for incompressible rules: the theorem covers all product hypotheses. For structured rules, error envelopes do & \cref{thm:teacher:main,prop:limits:envelope}\\
Is the three-hypothesis calculation realised by instances of the main theorem? & Yes. Its per-sample version is realised with an $11{,}710$-bit rule; its total-bits version (\cref{rem:teacher:hanni}) in population form, with $\lambda\in(9.77,58.0]$ at $n=10^4$, $\kappa=0$ (three hypotheses: $(10,100)$) & \cref{rem:teacher:hanni,rem:teacher:numbers,ex:teacher:version1}\\
Must the length coefficient grow with the data? & Yes: $\lambda=\Theta(n)$ is needed & \cref{prop:limits:plain,rem:limits:schedule}\\
Does an absolute bound on model complexity select the same hypotheses as the penalty? & In part: at vertices of the hull only & \cref{lem:setup:hull}\\
Is some other pathology present? & Yes: rate blindness, error envelopes, guard erosion, non-identifiability of $\lambda$, input-law dependence, labelled history, possible confident wrong extrapolation off the training support & \cref{sec:limits}, \cref{prop:absorb:history}\\
Do a sampler and its fill-ins get similar posterior weight at $\lambda=1$, and does a steeper prior prefer the sampler? & Yes, with explicit factors, in SFI and in SI. Yes for fixed $\lambda>1$ with private randomness (\cref{rem:split:reading}); no with shared randomness or in SI & \cref{thm:split:lambda1,prop:split:si,thm:split:private,thm:split:shared}\\
Against hypothesis length, does the entropy left given the hypothesis change slope at the model? & Yes for product models, at every flip rate below $1/4$; not for set models & \cref{prop:algstat:parity}\\
Has a steeper penalty been studied for the model--noise split? & In part: $\lambda$-weighted MDL; restricted model families & \cref{sec:related}\\
\bottomrule
\end{tabularx}
\caption{Questions studied in this paper and where they are answered. Numbers with $\kappa=0$ set the machine constant to $0$ and are illustrations.}\label{tab:intro:claims}
\end{table}

\paragraph{Scope and reader's guide.} The positive result is proved only for parities, error sets of product form, uniform inputs, flip rates below $0.056$ and incompressible error rules (\cref{rem:teacher:hypotheses}); for simple, realistic error patterns we have counterexamples but no positive result. For sampled data, near-population constants need about $10^{11}$ samples; coarser guarantees hold from about $10^9$ (\cref{cor:teacher:sampled}). With $\lambda=cn$ the learner is a family of batch objectives, not one Bayesian predictor (\cref{rem:setup:persample}). After related work (\cref{sec:related}) and the setting (\cref{sec:setup}), \cref{sec:absorb,sec:teacher,sec:limits,sec:split,sec:algstat,sec:computable} contain the results, each section with an appendix of proofs, and \cref{sec:discussion} collects the open problems. Notation follows the companion paper.
